\section{Conclusion}
\label{sec:conclusion}
% LOCKED: Validated May 15 conclusion wording. Keep this as a single paragraph.
% Requirements: plain-English takeaway, no future-work sentence, numeric evidence
% only where it proves the main claims, and explicit external-testbed comparison.

Our controlled threat-aware evaluation framework shows that measured quantum
entanglement-routing robustness varies with the joint interaction
of path selection, qubit allocation, replay capacity, and threat response. The
matched evaluation reveals a performance hierarchy within the context-aware neural
tier: \texttt{iCPursuitNeuralUCB} reaches 90.9\% internal average efficiency and
clears the 90\% model-family threshold, while the closest internal alternatives
remain below that mark. Across harder threat regimes, the framework further
exposes the fragility of adversarial-first exploration: \texttt{EXPNeuralUCB}
falls to an 18\% robustness floor, while the strongest contextual baselines
maintain floors near 77--81\%. The same controlled evaluation reveals a
threat-dependent capacity paradox, showing that replay capacity must be matched to
the active threat and that allocator choice directly shapes routing robustness.
Within the evaluated settings, the performance hierarchy and cross-threat
robustness results identify a candidate fixed routing baseline when the threat is
unknown; when the threat is detected, they support adapting allocator and replay
settings to the threat regime. The external-testbed comparison, which extends the
evaluation to four additional quantum-network testbeds, shows that the observed
performance hierarchy persists beyond the primary topology:
\texttt{iCPursuitNeuralUCB} leads average efficiency on \paperTwo{} (74.5\%,
95/300 wins), \paperSeven{} (78\%, 245/300 wins), \paperTwelve{} (44\%,
97/300 wins), and \paperEight{} (67.9\%). The advantage is clearest on the dense
50-node \paperSeven{} setting, where it reaches 78\% versus 70.8\% for the next
strongest models, and it still remains best on the larger \paperTwelve{} fusion
topology, where methods compress tightly at 44\% versus 43.8\%.

% -----------------------------------------------------------------------------
% Previous result-heavy one-paragraph conclusion retained for reference.
% -----------------------------------------------------------------------------
% This work evaluated quantum entanglement routing as a joint path-selection and
% qubit-allocation problem under matched threat, allocator, and replay-capacity
% settings. Across the validated corpus, pursuit-based contextual and neural
% hybrids, \texttt{CPursuitNeuralUCB} and \texttt{iCPursuitNeuralUCB}, produced
% the strongest deployment results: under the fixed 6K, \texttt{Fixed}-allocator,
% $T$-type, $s{=}2$ slice, they achieved 92.4--96.0\% Oracle-normalized
% efficiency with cross-scenario CVs of 6.5\% and 3.3\%, while contextual policies
% maintained 77--81\% robustness floors across the Markov, Adaptive, and
% OnlineAdaptive threat regimes and adversarial-first \texttt{EXPNeuralUCB} fell
% to an 18.0\% floor. The main result is that routing robustness is jointly
% determined by the learning rule, replay capacity, allocator choice, and threat
% regime: under \texttt{ThompsonSampling} allocation and $T_b$-type replay
% anchoring, OnlineAdaptive efficiency drops from 89.2\% to 84.9\% as replay scale
% increases from $s{=}1$ to $s{=}1.5$, then recovers to 90.8\% at $s{=}2$,
% revealing a threat-conditioned capacity paradox; changing only the allocator
% shifts the worst-case floor for the same \texttt{iCPursuitNeuralUCB} learner by
% 15.6~pp in the validated RQ3c 6K allocator comparison slice
% (\texttt{Fixed} floor~=~88.9\% vs.\ \texttt{Thompson} floor~=~73.3\%), making
% allocator policy part of the robustness mechanism. These findings support a
% concrete deployment rule: a static default of
% \texttt{iCPursuitNeuralUCB\,+\,Fixed\,+\,($T$-type, $s{=}2$)} sustains a
% 92.8\% worst-case floor and 96.0\% mean efficiency across all five threats at
% 6K, while allocator and replay settings should be adjusted to the detected
% threat regime. Cross-testbed results confirm the same trend while showing that
% larger or more constrained topologies narrow the Oracle gap.

% -----------------------------------------------------------------------------
% Previous three-paragraph conclusion retained for reference.
% -----------------------------------------------------------------------------
% This work evaluated quantum entanglement routing as a joint path-selection and qubit-allocation problem under matched threat, allocator, and replay-capacity settings. Across the validated corpus, the strongest deployment results came from pursuit-based contextual and neural hybrids, \texttt{CPursuitNeuralUCB} and \texttt{iCPursuitNeuralUCB}. Under the fixed 6K, \texttt{Fixed}-allocator, $T$-type, $s{=}2$ deployment slice, these models achieved 92.4--96.0\% Oracle-normalized efficiency with cross-scenario CVs of 6.5\% and 3.3\%, respectively. In the broader adversarial-scope RQ2 comparison, contextual policies maintained robustness floors of 77--81\%, while adversarial-first \texttt{EXPNeuralUCB} collapsed to an 18.0\% floor.
%
% The main result is that routing robustness is jointly determined by the learning rule, replay capacity, allocator choice, and threat regime. Under \texttt{ThompsonSampling} allocation and $T_b$-type replay anchoring, OnlineAdaptive efficiency drops from 89.2\% to 84.9\% as replay scale increases from $s{=}1$ to $s{=}1.5$, then recovers to 90.8\% at $s{=}2$. We refer to this threat-conditioned non-monotonicity as a capacity paradox. Similarly, changing only the allocator can shift the worst-case floor for the same \texttt{iCPursuitNeuralUCB} learner by 15.6~pp in the validated RQ3c 6K allocator-comparison slice (\texttt{Fixed} floor~=~88.9\% vs.\ \texttt{Thompson} floor~=~73.3\%), which means allocator policy is part of the robustness mechanism.
%
% These findings support a concrete deployment rule: a static default of \texttt{iCPursuitNeuralUCB\,+\,Fixed\,+\,($T$-type, $s{=}2$)} sustains a 92.8\% worst-case floor and 96.0\% mean efficiency across all five threats at the 6K horizon; allocator and replay settings should then be adjusted based on the detected threat regime. Our cross-testbed results confirm the same direction of effect, while also exposing scale limits as larger or more constrained topologies narrow the gap to the Oracle. Future work should expand the policy class beyond the current pursuit--neural hybrids, including models suggested by the strongest observed gains, and stress test the capacity paradox over longer horizons, larger replay scales, and more diverse threat and topology regimes.


% DK: Leave this in. Will need this ack section
\subsection*{Acknowledgements}
\noindent This material is based upon work supported by [hidden].

% SEIS: For all papers, references may extend as many pages beyond the page limit as you need.

%\noindent This material is based upon work supported by the [Hidden] under grant [Hidden]
